%!TEX TS-program = xelatex
%!TEX encoding = UTF-8 Unicode
\documentclass[a4paper, 12pt]{article}

% DEFAULT PACKAGES
\usepackage{ifxetex}
\ifxetex
  \usepackage{fontspec}
  \usepackage{xltxtra}
  \usepackage[czech]{babel}
\else
  \usepackage[cp1250]{inputenc}
  \usepackage[czech]{babel}
\fi
\usepackage{diplomova_prace}

% MY PACKAGES
\usepackage{makecell}
\usepackage{caption}
\usepackage{parskip} % just by calling this package: paragraph indentation will be suppressed, instead there will be a space after
\usepackage[margin=2.5cm]{geometry}
\usepackage{amsmath}

\usepackage{tabularx}
\newcolumntype{Z}{>{\raggedright\arraybackslash}X} % new column, which has responsive width like X, but is left-aligned
\renewcommand{\arraystretch}{1.2} % space after each row in table

\usepackage{chngcntr} % numbering for figures, tables and equations will include section
\counterwithin{figure}{section}
\counterwithin{table}{section}
\counterwithin{equation}{section}

% MY COMMANDS
\defineshorthand{"-}{\babelhyphen{nobreak}} % non-breakable hyphen
\def\D{\mathrm{d}} % non-italic differential sign

% BEGIN
\begin{document}
\pagenumbering{gobble} % don't number pages until after TOC

% souhrn
\section*{SOUHRN}
Souhrn v ČJ

\section*{SUMMARY}
Souhrn v AJ

% podekovani
\newpage ~ \vspace{17cm} \section*{PODĚKOVÁNÍ}
Rád bych poděkoval atd\ldots

% insert TOC, and start numbering pages
\newpage \tableofcontents \newpage
\pagenumbering{arabic}



\section{ÚVOD} \label{intro}
Kaprolaktam je významnou průmyslovou chemikálií, které se ročně vyrobí zhruba ?? tun.
V Evropě jsou jeho největší producenti ???, v České republice je to podnik Spolana s.r.o. (dále jen Spolana) v Neratovicích.
Jeho hlavní využití spočívá v polymeraci na polyamid, který se označuje jako Nylon 6 a uplatnění nalézá především na výrobu umělých vláken pro textilní průmysl i speciální aplikace, nebo jako pevný a houževnatý konstrukční materiál pro strojírenství <<???CITACE: Ullmann???>>.

Je známo mnoho postupů pro průmyslovou výrobu kaprolaktamu, tím nejběžnějším, který je využíván i ve Spolaně, je výroba z cyklohexanonu přes cyklohexanon-oxim pomocí kyseliny sírové.
Při výrobě vznikají dvě nemísitelné fáze -- organická, skládající se převážně z kaprolaktamu, a vodný roztok síranu amonného, který obsahuje zhruba 1~\% kaprolaktamu.
Tato práce se zabývá právě izolací zbytkového kaprolaktamu z vodné fáze, což se v současnosti provádí kapalinovou extrakcí trichlorethylenem, který je však v rámci evropské směrnice REACH zařazen mezi látky vzbuzující mimořádné obavy (SVHC) i na seznam látek podléhajících povolení (příloha XIV)~\cite{ECHA_TCE}, proto jsou hledány alternativy.

Tato práce poslouží k vypracování technologických podkladů pro náhradu trichlorethylenu toluenem, což se jeví jako velmi perspektivní.
V rámci práce byla tedy nejprve v laboratoři naměřena data popisující rovnováhy kapalina-kapalina v rozsahu, který je pro řešení problému praktický,
následně byl vypracován program v programovém prostředí MATLAB®, který obsahuje dva matematické modely, a to zaprvé model patrový, sloužící k vyhodnocení výkonu laboratorního či průmyslového extraktoru, zadruhé model přenosový, který pomocí základních veličin umožní vypočítat buďto návrh geometrie a provozních parametrů extraktoru, nebo simulaci definovaného zařízení pro předpověď jeho separační účinnosti.
Kromě toho byl sestaven laboratorní model extraktoru, aby mohla být funkce extraktoru testována s provozními vzorky.

Výsledky této práce budou využity k návrhu extraktorů síranových louhů v rámci nové technologie na separaci kaprolaktamu, která již nebude využívat chlorovaných rozpouštědel, a sníží se tak zátěž životního prostředí v České republice.



\newpage \section{LITERÁRNÍ ČÁST} \label{lit}
Literární část této práce se bude zabývat výčtem současně užívaných technologií na výrobu kaprolaktamu (dále jen CPL), možnostmi zařízení na průmyslovou kapalinovou extrakci, matematickým popisem extraktoru a teorií rovnováhy kapalina-kapalina.

\subsection{Technologie kaprolaktamu} \label{CPLtechnology}
\textit{tohle je málo - rozvést alternativy (subsubsection)}

$\varepsilon$-Kaprolaktam je bílá krystalická látka s bodem tání 69,0°C~\cite{CPLSpolanaList}.
Tato sloučeina se sumárním vzorcem \ce{C6H11NO} je z chemického hlediska laktam $\varepsilon$-aminokapronové kyseliny, molekula tedy tvoří sedmičetný cyklus.
V literatuře je popsáno mnoho technologií na výrobu kaprolaktamu, z toho nejvýznamnější (?? \% světové produkce) je však postup popsaný v kapitole \ref{CHON}, který je zaveden i ve Spolaně a týká se jej tato práce.

\subsubsection{Tradiční výroba z cyklohexanonu} \label{CHON}
\textit{udělat reakční schéma, napsat proč je obsah síranu takový či makový}

Surovinou pro tuto technologie je cyklohexanon, typicky vyráběný oxidací cyklohexanu. Ten je nejprve substituován hydroxylaminem na cyklohexanon-oxim <<???CITACE: Ullmann???>> a následně podroben Beckmannově přesmyku katalyzovaném koncentrovanou kyselinou sírovou.
Vzniklá sloučenina kyseliny sírové a kaprolaktamu je nakonec neutralizována roztokem amoniaku, v důsledku čehož se reakční směs rozdělí na dvě fáze:
organickou, která obsahuje většinu kaprolaktamu, vedlejších produktů i část vody, a vodný roztok síranu amonného (dále jen SA) o koncentraci okolo 40~\%, který však vedle stopových nečistot obsahuje také zhruba 1~\% hm. CPL (v této práci bude složení uváděno jako hmotnostní zlomek, pokud nebude uvedeno jinak).

\subsection{Metody extrakce} \label{extraction}
\textit{stručný popis protiproudé extrakce, jaké jsou konstrukční možnosti, KARR UPSCALING}

\subsection{Matematický popis extraktoru} \label{the_Math}
Kontinuální protiproudý extraktor v ustáleném stavu lze matematicky popsat buďto patrovým modelem nebo přenosovým modelem, tyto modely se však zásadně liší požadovanými informacemi i možností použití.
Patrový model, který bude dále popsán až v kapitole~\ref{compStage}, má využití zejména na porovnání separační účinnosti reálných extraktorů na základě změřených bilančních veličin, v případě Karrova extraktoru i na zvětšování měřítka z laboratorních experimentů na průmyslové zařízení, jak bylo popsáno v kapitole~\ref{extraction}.

Naproti tomu přenosový model k problému přistupuje jiným způsobem -- cílem je namodelovat chování extraktoru pouze ze znalostí základních vlastností, jedná se o vlastnosti látek, jako je hustota, dynamická viskozita a difuzivita kaprolaktamu v dané látce, a také o vlastnosti celého systému, jako je mezifázové napětí a popis rovnováhy kapalina-kapalina v podobě rovnovážné funkce: $u=f(w)$.

Ve tomto textu budou dále používány m.j. tyto veličiny: $t$ znamená teplota v °C, $\dot{m}$ je hmotnostní tok, $Q$ je objemový tok, $u$ a $w$ jsou obecně hmotnostní zlomky CPL v organické fázi, resp. ve vodné fázi.
Veličiny s dolními indexy F, R, S a E se vztahují popořadě k surovině, rafinátu, solventu a extraktu.

\subsubsection{Přenosový model s axiální disperzí} \label{diff2}
V této kapitole bude podrobně vysvětlen výpočetní postup, který popisuje extraktor jako soustavu dvou nelineárních diferenciálních rovnic druhého řádu, kde každá rovnice popisuje jednu fázi.
Nezávislou veličinou modelu je bezrozměrná výška $z$ definovaná jako podíl výškové souřadnice $l$ vůči celkové výšce extraktoru $L$, a to ve směru od suroviny k rafinátu (či od extraktu k solventu).
Závislé veličiny jsou $c_x$ a $c_y$, tedy hmotnostní koncentrace CPL ve fázi ochuzované, resp. obohacované o CPL, což v případě této práce vždy znamená fázi vodnou, resp. organickou.
Tento výpočetní postup, který popsali S.~Hartland a A.~Kumar~\cite{HartlandKumar1997}, je vyjádřen soustavou rovnic~\ref{eqn:diff2}, kde $Pe_x$ a $Pe_y$ jsou Pécletova čísla (definovaná v rovnicích~\ref{eqn:Peclet} a~\ref{eqn:Peclet2}), $T_x$ je počet přenosových jednotek vzhledem k fázi x (definovaný v rovnici~\ref{eqn:Tx}) a $c_y^*$ je taková hodnota $c_x$, která by byla v rovnováze s $c_y$.
\begin{align}
\label{eqn:diff2}
	\frac{\D^2 c_x}{\D z^2} - Pe_x \frac{\D c_x}{\D z} &=  Pe_x T_x (c_x - c_y^*) \\
	\frac{\D^2 c_y}{\D z^2} + Pe_y \frac{\D c_y}{\D z} &= -Pe_y T_x (c_x - c_y^*) \frac{Q_x}{Q_y} \nonumber
\end{align}
\begin{align}
\label{eqn:Peclet}
	Pe_x = \frac{Q_x L}{A E_x} \\
\label{eqn:Peclet2}
	Pe_y = \frac{Q_y L}{A E_y}
\end{align}
$E_x$ a $E_y$ v rovnicích ~\ref{eqn:Peclet} a~\ref{eqn:Peclet2} jsou koeficienty axiální disperze.
$K_x$ v rovnici~\ref{eqn:Tx} je koeficient prostupu hmoty vztažený k fázi x, $a$ je specifická mezifázová plocha a $A$ je plocha průřezu extraktoru.
\begin{equation}
	T_x = \frac{K_x a A L}{Q_x}
\label{eqn:Tx}
\end{equation}

Řešení soustavy rovnic~\ref{eqn:diff2} musí splňovat okrajové podmínky, které jsou pro bod $z=0$ definovány v rovnicích~\ref{eqn:bound2_1} a~\ref{eqn:bound2_2}.
\begin{align}
\label{eqn:bound2_1}
	\left( \frac{\D c_x}{\D z} \right)_{z=0} &= Pe_x (c_{x,0} - c_{x,\mathrm{in}}) \\
\label{eqn:bound2_2}
	\left( \frac{\D c_y}{\D z} \right)_{z=0} &= 0
\end{align}
V bodě $z=1$ musí být splněny podmínky definované v rovnicích~\ref{eqn:bound2_3} a~\ref{eqn:bound2_4}.
\begin{align}
\label{eqn:bound2_3}
	\left( \frac{\D c_x}{\D z} \right)_{z=1} &= 0 \\
\label{eqn:bound2_4}
	\left( \frac{\D c_y}{\D z} \right)_{z=1} &= Pe_y (c_{y,\mathrm{in}} - c_{y,1})
\end{align}
Řešením této soustavy pak je funkce $(c_x, c_y) = f(z)$, přičemž je nutné podotknout, že krajní hodnoty $c_x$ a $c_y$ odpovídají výstupním proudům, avšak nikoliv vstupním proudům, tedy $c_x(0) \neq c_\mathrm{F}$, $c_y(0) = c_\mathrm{E}$, $c_x(1) = c_\mathrm{R}$ a $c_y(1) \neq c_\mathrm{S}$. Při vstupech proudů dochází ke skokovým změnám koncentrací, lze tedy říci, že s vyšší hodnotou koeficientů axiální disperze se charakter toku obou fází vzdaluje pístovému toku a přibližuje ideálnímu mísiči.

Při výše uvedeném výpočtu je $L$ považováno za konstantu, takovýto výpočet tedy slouží jako výpočet simulační, zrovna tak jej však lze realizovat jako výpočet návrhový, kdy $L$ je v soustavě dif. rovnic neznámým parametrem, který je hledán na základě počátečního odhadu.
Při zavedení neznámého parametru je však nutné specifikovat další okrajovou podmínku, kterou je deklarace požadované koncentrace rafinátu (rovnice~\ref{eqn:boundExtra}).
\begin{equation}
	\left( c_x \right)_{z=1} = c_\mathrm{R}
\label{eqn:boundExtra}
\end{equation}

Programová implementace tohoto modelu je popsána v kapitole~\ref{compTransfer}.

\subsubsection{Přenosový model s pístovým tokem} \label{diff1}
Popsaný matematický model lze zjednodušit pomocí předpokladu pístového toku, pak jsou koeficienty axiální disperze nulové a lze odvodit soustavu rovnic~\ref{eqn:diff1}.
\begin{align}
\label{eqn:diff1}
	\frac{\D c_x}{\D z} &= -T_x (c_x - c_y^*)\\
	\frac{\D c_y}{\D z} &= -T_x (c_x - c_y^*) \frac{Q_x}{Q_y} \nonumber
\end{align}
Při pístovém toku jsou však odlišné okrajové podmínky -- nedochází ke skokovým změnám koncentrací, okrajové podmínky jsou tedy vlastně deklaracemi krajních hodnot $c_x$ a $c_y$, jak je uvedeno v rovnicích~\ref{eqn:bound1_1}, ~\ref{eqn:bound1_2}.
Při návrhovém výpočtu je zde, obdobně jako u modelu s axiální disperzí, jedna okrajová podmínka navíc, která již byla definována v rovnici~\ref{eqn:boundExtra}.
\begin{align}
\label{eqn:bound1_1}
	\left( c_x \right)_{z=0} &= c_\mathrm{F} \\
\label{eqn:bound1_2}
	\left( c_y \right)_{z=1} &= c_\mathrm{S}
\end{align}

Model s pístovým tokem je využit jako mezivýpočet při implementaci axiálně disperzního modelu, jak je popsáno v kapitole~\ref{compTransfer}.

\subsection{Rovnováha kapalina-kapalina} \label{LLEtheory}
\textit{Krátký teoretický popis, proč my na teorii kašleme, nějaká data z literatury, ta ale nestačí a proč}



\newpage \section{EXPERIMENTÁLNÍ ČÁST} \label{exp}

\subsection{Použité chemikálie a materiály} \label{materials}
Tabulka \ref{tab:chemicals} je výčtem všech použitých chemikálií.

\begin{table}[ht]
\caption{Seznam použitých chemikálií} \label{tab:chemicals}
\begin{tabularx}{\textwidth}{llp{2cm}Z}
	\Xhline{2pt} \textbf{látka} & \textbf{výrobce} & \textbf{čistota} & \textbf{H/P věty} \\ \hline
	destilovaná voda & VŠCHT & & \\
	propan-2-ol & Penta & \makecell{p.a.\\>99,8~\%} & H225, H319, H336, P210, P233, P305+P351+P338 \\
	n"-heptanol & Fluka & \makecell{p.a.\\>99,0~\%} & H319, P305+P351+P338 \\
	n-heptan & Penta & \makecell{p.a.\\>99,0~\%} & H225, H304, H315, H336, H410, P210, P243, P273, P301+P310, P331 \\
	toluen & Penta & \makecell{p.a.\\>99,0~\%} & H225, H304, H315, H336, H361d, H373, P210, P261, P280, P301+P310, P331 \\
	trichlorethylen & Spolana & \makecell{destilovaný\\99,7~\%} & H315, H317, H319, H341, H350, H412, P201, P273, P280, P302+P352, P305+P351+P338, P308+P313 \\
	kaprolaktam & Spolana & \makecell{destilovaný\\100,0~\%} & H319, H315, H332, H302, H335, P261, P280, P304+P340, P305+P351+P338, P301+P330, P403+P233 \\
	síran amonný & Lachema & \makecell{p.a.\\>99,0~\%} & H315, H319, H335, P280, P305+P351+P338 \\
	\Xhline{2pt}
\end{tabularx}
\end{table}

\subsection{Měření rovnováhy kapalina-kapalina} \label{LLE}
\textit{lorem ipsum}

\subsubsection{Aparatury a metodika experimentů} \label{LLE_exp}
K provádění experimentů na měření dat LLE byla použita aparatura znázorněná na obrázku \ref{fig:LLE_a}. Aparaturu tvoří skleněná cela, ve které je připravena fáze těžší~(1) i lehčí~(2).
Jejich disperzi zajišťuje magnetické míchadlo~(3), k udržení konstantní teploty slouží duplikátor, do nějž vstupuje~(4) a vystupuje~(5) voda z termostatu.
Těchto cel bylo k~dispozici celkem sedm. V~celách byly připraveny příslušné kapalné směsi, a~to~tak, aby objem obou fází činil zhruba 20~ml, načež bylo nastaveno míchání dostatečně intenzivní, aby vznikla jemná disperze kapalných fází.
Po zapnutí termostatu s příslušným teplotním nastavením bylo vyčkáno na ustálení teploty, poté bylo míchání ponecháno v běhu minimálně 30~min, aby se systém dostal dostatečně blízko k fázové rovnováze.
Po uplynutí této doby bylo míchání vypnuto a systém byl ponechán potřebnou dobu v klidu, aby došlo k úplnému mechanickému oddělení kapalných fází, tedy aby byly obě fáze vizuálně čiré (minimálně však 5~min).
Nakonec byly z obou fází odebrány vzorky o hmotnosti cca. 1~g, pro účely analýzy (kapitola \ref{GC}) byla vždy změřena také jejich přesná navážka na analytických vahách.

\begin{figure}[ht]
\begin{center}
	\includegraphics[width=6cm,height=6cm,keepaspectratio,]{../../images/LLE_cela.png}
	\caption{Aparatura na měření LLE} \label{fig:LLE_a}
\end{center}
\end{figure}

\subsubsection{Plynová chromatografie} \label{GC}
Vzorky organické i vodné fáze jak z měření dat LLE, tak z experimentů s laboratorním extraktorem (kapitola \ref{labmodel}), bylo třeba podrobit kvantitativní analýze, k čemuž byl využit plynový chromatograf (GC) Shimadzu GC-2014 vybavený plamenově ionizačním detektorem (FID).
Měřen byl pouze obsah CPL, neboť vzhledem k rozsahu měřených koncentrací lze předvídat nemísitelnost organických a vodných fází, i obsah nečistot se ukázal být pouze stopový.
Na měření obsahu SA ve vodné fázi nebyly k dispozici prostředky, takže ten byl zajištěn pouze navážkou při přípravě roztoku.
Pro účely analýzy tedy bylo předpokládáno, že vzorky organické fáze obsahují pouze CPL a solvent (TCE či TOL), zatímco vzorky kapalné fáze obsahují CPL, vodu a SA.

Stanovení obsahu CPL bylo provedeno metodou vnitřního standardu, kterým byl n"-heptanol -- do každého vzorku o známé navážce bylo na analytických vahách odváženo zhruba dvojnásobné množství standardního roztoku, který obsahuje přibližně 0,5~\% n"-heptanolu v isopropylalkoholu.
Tento roztok plní nejen funkci vnitřního standardu, ale také jako srážecí činidlo pro SA -- vzorky vodné fáze s rozpuštěnou solí by nebylo možné analyzovat, neboť by mohlo dojít k ucpání chromatografické kolony.
Po vysrážení soli tedy byla odebrána pouze kapalná fáze, která byla před nástřikem na GC ještě odstředěna na laboratorní centrifuze.

Parametry kolony a podmínky při chromatografii jsou uvedeny v tabulce \ref{tab:GC}.
Teplotní profil byl zvolen tak, aby byl retenční čas kaprolaktamu (nejméně těkavá látka z analytů) co nejnižší, a tedy aby jeho pík měl tvar co nejbližší Gaussově funkci, ale zároveň aby byly píky kaprolaktamu i n"-heptanolu dostatečně oddělené od sebe i od píků rozpouštědel.
K tomuto účelu byl přizpůsoben následující teplotní profil: 100°C po dobu 2~min, ohřev na teplotu 200°C trvající 4~min a nakonec 2~min při této teplotě.
Tabulka \ref{tab:GCtime} shrnuje přibližné hodnoty retenčních časů ($\tau_\mathrm{R}$), které byly za těchto podmínek pozorovány.
V případě, že je retenční čas zadán jako rozmezí, se jedná o pík s šikmým tvarem a jeho vrchol se proto pohybuje v závislosti na celkové ploše píku.

\begin{table}[ht]
\caption{Podmínky při analýze GC a popis kolony} \label{tab:GC}
\begin{center}
\begin{tabular}{ll}
	\Xhline{2pt}
	chromatograf & Shimadzu GC-2014\\
	ovládací software & Clarity 7.0 (DataApex)\\
	kolona & SPB-5 (Výrobce Supelco, Inc.)\\
	délka kolony & 30 m\\
	průměr kolony & 0,32 mm\\
	tloušťka filmu & 0,25 \textmu m\\
	splitovací poměr & 15\\
	objem nástřiku & 0,6 \textmu l\\ % ??? nejsem si jist, zda-li je to v mikrolitrech
	teplota injektoru & 250°C\\
	teplota detektoru & 250°C\\
	nosný plyn & He\\
	přetlak na vstupu do kolony & 72 kPa\\
	střední rychlost nosného plynu & 33,0 cm/s\\
	\Xhline{2pt}
\end{tabular}
\end{center}
\end{table}

\begin{table}[ht]
\caption{Retenční časy analytů na GC} \label{tab:GCtime}
\begin{center}
\begin{tabular}{ll}
	\Xhline{2pt}
	\textbf{Analyt} & \textbf{$\tau_\mathrm{R}$/min}\\ \hline
	propan-2-ol & 1,80\\
	toluen & 2,17\\
	n"-heptanol & 3,01\\
	kaprolaktam & 5,17--5,46\\
	\Xhline{2pt}
\end{tabular}
\end{center}
\end{table}

\subsubsection*{Zpracování chromatogramů}

Po analýze vzorku v plynovém chromatografu poskytne ovládací software chromatogram a z něj vyhodnotí různé veličiny, z nichž nejdůležitější je retenční čas, který slouží pro přiřazení píku určité látce, a plocha píku, pomocí níž lze obsah látky kvantifikovat.

Při využití metody vnitřní kalibrace lze využít vzorec \ref{eqn:GC_cal}, kde $w_\mathrm{CPL,s}$ a $w_\mathrm{Hep,s}$ značí obsah CPL a n"-heptanolu v analyzovaném vzorku (po přidání standardu), $A_\mathrm{CPL}$ a $A_\mathrm{Hep}$ značí plochy píků a $S_{^\mathrm{CPL}/_\mathrm{Hep}}$ je kalibrační koeficient kaprolaktamu vůči n"-heptanolu.
Při známých navážkách původního vzorku i standardu je tedy nejprve vypočten $w_\mathrm{Hep,s}$, následně $w_\mathrm{CPL,s}$ a nakonec lze určit $w_\mathrm{CPL}$, tedy obsah CPL v původním vzorku.

\begin{equation}
	w_\mathrm{CPL,s} = S_{^\mathrm{CPL}/_\mathrm{Hep}} w_\mathrm{Hep,s} \frac{A_\mathrm{CPL}}{A_\mathrm{Hep}}
\label{eqn:GC_cal}
\end{equation}

Pro zpracování chromatogramů však bylo nejprve nutné provést kalibraci pomocí roztoků se známým obsahem n"-heptanolu, který je při dané metodice víceméně konstantní, i kaprolaktamu, jehož obsah se může mezi měřenými vzorky řádově lišit.
Kalibrační vzorky proto byly připraveny v širokém rozsahu koncentrací, a to od 0.004--10.2~\% včetně standardu, což odpovídá 0.012--30.4~\% ve vzorku bez standardu.
Pro každý kalibrační vzorek byla určena hodnota $S_{^\mathrm{CPL}/_\mathrm{Hep}}$, a jelikož koeficient nevykazoval trend vůči obsahu CPL, byla v této práci používána průměrná hodnota.

\subsection{Měření mezifázového napětí} \label{gamma}
Pro účely sestavení matematických modelů se později ukázalo, že je třeba změřit také mezifázové napětí, neboť sice lze v literatuře nalézt vzorce na jeho odhad~\cite{Demond1993,Perry_15}, ty však bývají velmi nespolehlivé.
Proto bylo provedeno měření pro systém \ce{H2O}--SA--CPL--TOL, aby mohla být zjištěna závislost mezifázového napětí na koncentraci CPL ve vodné fázi.

Měření bylo provedeno pomocí kapkové metody (též známá jako stalagmometrická metoda), která je klasicky využívána pro měření povrchového napětí, lze ji však modifikovat pro měření mezifázového napětí, a to tak, že je nádoba stalagmometru naplněna vzorkem těžší fáze a kapilára je ponořena pod hladinu vzorku lehčí fáze, kde se tak tvoří rozhraní obou fází.
Postup měření je takový, že se nastaví takový vhodný tok těžší fáze, aby vznikaly kapky dostatečně pomalu, a ty jsou pak váženy v sériích čítajících alespoň dvacet kapek.
Doba mezi odtržením kapek musí být dostatečně dlouhá, jinak je měření zkresleno jevem, kdy přitékající kapalina kapku zvětší ještě při jejím odtrhávání, což vede k vyšší pozorované hmotnosti~\cite{Harkins1919}.

Mezifázové napětí~($\gamma$) lze vypočítat postupem, který popsali J.~L.~Lando a H.~T.~Oakley~\cite{Lando1967}. Vztah je uveden v rovnici~\ref{eqn:gamma}, kde $\rho_\mathrm{H}$ a $\rho_\mathrm{L}$ znamená hustotu těžší, resp. lehčí fáze, $m$ je hmotnost jedné kapky, $g$ tíhové zrychlení, $r$ poloměr kapiláry a $\Phi$ je korekční faktor, pro který autoři uvádějí empirický vztah. Ten je zde uveden v rovnici~\ref{eqn:gammaCorr}.
\begin{equation}
	\gamma = \frac{\rho_\mathrm{H} - \rho_\mathrm{L}}{\rho_\mathrm{H}} \frac{m g \Phi}{2 \pi r}
\label{eqn:gamma}
\end{equation}
\begin{equation}
	\Phi = -1.044 \left(\frac{r}{\sqrt[3]{m / \rho_\mathrm{H}}}\right)^2 + 1.754 \frac{r}{\sqrt[3]{m / \rho_\mathrm{H}}} + 0.929
\label{eqn:gammaCorr}
\end{equation}
\subsection{Laboratorní model extraktoru} \label{labmodel}
\textit{lorem ipsum}

\subsubsection{Popis aparatury} \label{labmodel_a}
\textit{lorem ipsum}

\subsubsection{Metodika experimentů} \label{labmodel_m}
o analýze jen viz kapitolu \ref{GC}



\newpage \section{VÝPOČETNÍ ČÁST} \label{comp}
Stěžejní částí této práce je komplexní počítačový program vytvořený v programovém prostředí MATLAB, který se skládá se z několika oddělených, avšak na sebe navazujících částí, kterými jsou: zpracování dat LLE na rovnovážné funkce, patrový model extraktoru a model přenosový.
V následujících kapitolách budou popsány principy fungování programu a způsob výpočtů, zatímco samotný zdrojový kód a podrobná dokumentace jeho použití se nachází v příloze~\ref{sourceCode}.

\subsection{Zpracování rovnovážných dat} \label{compLLE}
Tato část programu slouží ke zpracování naměřených dat LLE na funkce popisující rovnováhu ve tvaru $u = f(w, t)$.
Vzhledem k tomu, že data byla měřena pouze v malém koncentračním rozmezí ($w$ nepřesahuje 1~\%), bylo možné použít empirickou rovnici namísto modelů aktivitních koeficientů, které, byť jsou termodynamicky konzistentní, nebylo možné spolehlivě proložit.
V daném rozsahu $w$ a $t$ plně postačuje funkce, která je polynomem třetího řádu vůči $w$ (avšak bez absolutního členu, neboť funkce musí procházet bodem $(0,0)$) a polynom prvního řádu (lineární funkce) vůči $t$, jak je popsáno v rovnici~\ref{eqn:eqFun0}. Tato funkce má 5 parametrů $p_1$--$p_5$.
\begin{equation}
	u = (p_1 w^3 + p_2 w^2 + p_3 w)(p_4 t + p_5)
\label{eqn:eqFun0}
\end{equation}
V tomto tvaru jsou však parametry na sobě jasně závislé.
Numerický algoritmus na nelineární regresi, který je v prostředí MATLAB poskytován, potom hledá optimum obtížně a nespolehlivě, proto byla funkce použita v jiném tvaru, který je popsán v rovnici~\ref{eqn:eqFun}. Tato funkce má 6 parametrů $p_1$--$p_6$.
\begin{equation}
	u = p_1 w^3 t + p_2 w^2 t + p_3 w t + p_4 w^3 + p_5 w^2 + p_6 w
\label{eqn:eqFun}
\end{equation}
Do regrese by sice bylo možné zahrnout obsah síranu amonného ve vodné fázi jako třetí nezávislou proměnnou, regrese by tím však byla méně spolehlivá.
Obsah SA lze přitom považovat za zaručený, z důvodů popsaných v kapitole~\ref{CHON} jej navíc ani není možné příliš měnit, proto byly jako nezávislé proměnné zahrnuty pouze $w$ a $t$.

Výše popsaná nelineární regrese je implementována ve funkci \textit{fitEqCurves.m}, kde je využita vestavěná funkce \textit{lsqcurvefit} nabízená prostředím MATLAB, aby byla získána sada parametrů pro každou dvojici fází (vodná fáze s daným obsahem SA a vybrané rozpouštědlo).
Po úspěšném proložení regresních křivek lze spouštět funkci \textit{getEqFun.m}, která do rovnovážné funkce dosadí zadanou teplotu a získanou funkci $u = f(w)$ tabeluje jako vektory $u$ a $w$.
Tabelace funkce je výhodná pro výpočetně nenáročnou implementaci funkce inverzní, tedy $w = f^{-1}(u)$, která je hojně využívána v modelech extraktorů.

Pro vizualizaci naměřených dat byla vytvořena ještě funkce \textit{plotLLEdata.m}, která do grafu závislosti $u(w)$ vykreslí pro zadanou dvojici fází všechna data LLE, která jsou seskupena do sad dle teploty při měření.

\subsection{Implementace patrového modelu} \label{compStage}
Rovnovážné funkce využívá patrový model extraktoru, který ze zadané bilance vypočte počet teoretických pater (TP) pomocí McCabe"-Thielovy metody~\cite{McCabe1925} modifikované pro popis kapalinové extrakce, jak ji popisuje Perry's Chemical Engineers' Handbook~\cite{Perry_15} (strana~15"-45).
Hlavním účelem patrového modelu je vyčíslení separační účinnosti laboratorního či průmyslového extraktoru pomocí počtu TP za zadané hmotnostní bilance.
Kromě toho je také vypočtena minimální hodnota hmotnostního toku solventu ($\dot{m}_\mathrm{S,min}$), což je klíčová informace pro navrhování procesu, neboť počet TP potřebných pro separaci se zvyšuje s klesajícím tokem solventu a při hodnotě $\dot{m}_\mathrm{S,min}$ limitně dosahuje nekonečna, při nižších hodnotách pak separace není možná.
Centrální funkcí tohoto modelu je \textit{extractor.m}, kde se nejprve dopočte hmotnostní bilance a zkontroluje její smysluplnost, načež je vypočtena hodnota $\dot{m}_\mathrm{S,min}$ a nakonec je modifikovanou McCabe"-Thielovou metodou určen počet TP.

Algoritmus na výpočet $\dot{m}_\mathrm{S,min}$ funguje obdobně jako dobře známá metoda půlení intervalu, nejprve jsou tedy zvoleny dvě hodnoty toku solventu -- velmi nízká hodnota $\dot{m}_1$ a velmi vysoká hodnota $\dot{m}_2$, přičemž se očekává, že hodnota $\dot{m}_\mathrm{S,min}$ bude ležet mezi těmito krajními hodnotami.
Následuje cyklus, kdy je při každé iteraci určena hodnota $\dot{m}_3$ jako aritmetický průměr hodnot $\dot{m}_1$ a $\dot{m}_2$ a následně je zjištěno, zda-li má pracovní přímka průsečík s rovnovážnou křivkou pro hodnoty $\dot{m}_1$, $\dot{m}_2$ a $\dot{m}_3$.
Pokud je průsečík nalezen pro $\dot{m}_1$ a $\dot{m}_3$, bude v další iteraci platit $\dot{m}_1 = \dot{m}_3$, zatímco pokud je průsečík nalezen pouze pro $\dot{m}_1$, bude platit $\dot{m}_2 = \dot{m}_3$. Takovýchto iterací je provedeno 30, což zaručuje vysokou přesnost určení hodnoty $\dot{m}_\mathrm{S,min}$.

Protože tato funkce poskytuje jen strojový výstup, byla vytvořena také funkce \textit{extractorPlot.m}, která přehledně vypíše důležité výsledky do konzole prostředí MATLAB a vykreslí graf závislosti $u(w)$, kde je vynesena rovnovážná křivka, pracovní úsečka a úsečky znázorňující jednotlivá teoretická patra.
Pro usnadnění parametrické studie byla vytvořena ještě funkce \textit{extractorFeed.m}, která umožní vybranou vstupní veličinu měnit v definovaném rozmezí, výsledkem je pak tabulka závislosti počtu TP a $\dot{m}_\mathrm{S,min}$ na vybraném parametru.

\subsection{Implementace přenosového modelu} \label{compTransfer}
Jak již bylo dříve zmíněno, přenosový model při předpovědi chování extraktoru vychází ze znalostí základních vlastností látek či systému.
Relevantní vlastnosti látek jsou hustota ($\rho$), dynamická viskozita ($\mu$) a difuzivita kaprolaktamu v dané látce ($D$), vlastnosti systému jsou rovnovážná funkce $u = f(w)$ a mezifázové napětí ($\gamma$).
Vlastnosti látek byly exportovány z databáze simulačního programu Aspen Plus v širokém rozmezí teplot, aby je bylo možné dle potřeby interpolovat a získat tak hodnoty za libovolně zvolené teploty.
Hodnoty pochází z různých literárních zdrojů, a to <<<???CITACE???>>>.
Rovnovážná funkce tabelovaná jako vektor $u$ a $w$ je získána spuštěním funkce \textit{getEqFun.m} popsané v kapitole~\ref{compLLE}.
Naměřené hodnoty mezifázového napětí byly zpracovány na funkci $\gamma(w)$, k čemuž postačuje regrese naměřených dat pomocí polynomu druhého řádu (viz kapitolu~\ref{resgamma}).

Přenosový model může být realizován buďto jako výpočet návrhový, kdy jsou vypočteny potřebné rozměry zařízení na dosažení požadované hm. bilance, a nebo jako výpočet simulační, kdy jsou rozměry zařízení specifikovány a model předpovídá separační účinnost extraktoru, čili doplní hm. bilanci.

\subsubsection{Řešení diferenciálních rovnic} \label{compDesign}
Jádrem výpočtu je matematický model, který popisuje extraktor jako soustavu dvou nelineárních diferenciálních rovnic druhého řádu, jak bylo podrobně popsáno v kapitole~\ref{the_Math}.

S.~Hartland a A.~Kumar~\cite{HartlandKumar1997} uvádějí analytické řešení diferenciálních rovnic při návrhovém výpočtu -- řešením modelu s pístovým tokem je přímé vyjádření $L$, zatímco řešením modelu s axiální disperzí je rovnice transcendentní pro $L$, z níž lze $L$ určit numerickým algoritmem na hledáním kořene.
Analytická řešení ovšem platí pouze za předpokladu lineární rovnováhy, což, jak bude ukázáno v kapitole~\ref{resLLE}, v našem případě neplatí, rovnovážná funkce je totiž značně nelineární.
Autoři uvádí, že při nelineární rovnováze diferenciální rovnice nemají analytické řešení a je nutné je řešit numericky.

K numerickému řešení diferenciálních rovnic s okrajovou podmínkou MATLAB nabízí funkci \textit{bvp5c}, ta pro svůj běh potřebuje zaprvé funkci popisující diferenciální model prvního řádu, která ze závislých a nezávislých proměnných vypočítá derivace závislých proměnných, zadruhé funkci definující počáteční podmínky jakožto rovnice v anulovaném tvaru, zatřetí pak funkci sloužící jako počáteční odhad řešení.
V případě návrhového výpočtu je třeba zahrnout i počáteční odhad výšky extraktoru.

Protože algoritmy řešení okrajových problémů nebývají zcela robustní a řešení může často divergovat, je třeba poskytnout dobrý počáteční odhad řešení.
Model s axiální disperzí se bez kvalitního odhadu ukázal být velmi nestabilní, proto je nejprve numericky vypočten model s pístovým tokem, jehož řešení je snazší a zcela přitom postačuje jednoduchý odhad lineárních profilů koncentrace.
V případě návrhového výpočtu, kdy jsou koncentrace vstupních i výstupních proudů známy, jsou lineární profily od $c_F$ do $c_R$, resp. od $c_E$ do $c_S$, v případě simulačního výpočtu lze namísto $c_R$ a $c_S$ dosadit nulovou koncentraci.
Pro odhad výšky, je-li třeba, postačuje použít výšku ekvivalentu teoretického patra 0,5~m, která se vynásobí počtem TP získaným zavoláním patrového modelu.
Diferenciální rovnice prvního řádu i jejich počáteční podmínky byly již popsány v rovnicích~\ref{eqn:diff1}, \ref{eqn:bound1_1}, \ref{eqn:bound1_2} a~\ref{eqn:boundExtra}.

Nalezené řešení pro model s pístovým tokem má využití pro sestavení počátečního odhadu pro model s axiální disperzí, ten je však nutné nejprve upravit ze soustavy dvou dif. rovnic druhého řádu, jak jsou uvedeny v rovnici~\ref{eqn:diff2}, na soustavu čtyř rovnic prvního řádu.
To je možné provést zavedením derivací koncentrací jako dvou dalších závislých proměnných: $d_x = \D c_x / \D z$ a $d_y = \D c_y / \D z$.
Diferenciální model připravený pro numerický řešič je rozveden v rovnici~\ref{eqn:diff24}.
\begin{align}
\label{eqn:diff24}
	\frac{\D c_x}{\D z} &= d_x \\
	\frac{\D c_y}{\D z} &= d_y \\ \nonumber
	\frac{\D d_x}{\D z} &=  Pe_x d_x + Pe_x T_x (c_x - c_y^*) \\ \nonumber
	\frac{\D d_y}{\D z} &= -Pe_y d_y - Pe_y T_x (c_x - c_y^*) \frac{Q_x}{Q_y} \nonumber
\end{align}
Protože řešení modelu s pístovým tokem je ve formě vektorů hodnot $c_x$ a $c_y$ odpovídajících vektoru $z$, odhadovou funkci pro model s axiální disperzí lze realizovat jako algoritmus, který pro požadované $z$ nalezne dva sousedící body existujícího řešení, mezi nimiž požadované $z$ leží.
Hodnoty $c_x$ a $c_y$ jsou lineárně interpolovány, zatímco hodnoty $d_x$ a $d_y$ jsou vypočteny běžnou dvoubodovou formulí pro numerickou derivaci.
Při návrhovém výpočtu je součástí řešení modelu s pístovým tokem také vypočtená výška extraktoru -- pro počáteční odhad výšky modelu s axiální disperzí je použit její dvojnásobek.

\subsubsection{Kvantifikace parametrů přestupu hmoty} \label{compInsfun}
\textit{lorem ipsum}



\newpage \section{VÝSLEDKY A DISKUSE} \label{res}
\textit{lorem ipsum}

\subsection{Výsledky experimentů}
\textit{lorem ipsum}

\subsubsection{Rozsah měření rovnovážných dat}
\textit{lorem ipsum}

\subsubsection{Rovnovážná data} \label{resLLE}
\textit{lorem ipsum}

\subsubsection{Mezifázové napětí} \label{resgamma}
\textit{doba mezi kapkama nebyla dost dlouhá...}

\subsubsection{Laboratorní model extraktoru}
\textit{lorem ipsum}

\subsection{Výsledky modelů}
\textit{Nejdříve tabulka jednoho modelového zadání programu se verb 2}

\subsubsection{Parametrická studie}
\textit{Tabulky z callTable}

\subsubsection{Možnosti návrhu extraktorů}
\textit{lorem ipsum}

\subsubsection{Srovnání s laboratorním extraktorem}
\textit{lorem ipsum}



\newpage \section{ZÁVĚR} \label{zaver}
\textit{lorem ipsum}



\newpage \section*{LITERATURA}
\addcontentsline{toc}{section}{LITERATURA} % manually add TOC entry cause section* doesn't by default
\renewcommand\refname{\vskip -1cm} % clear the automatic title and remove empty space
\bibliography{references}

\newpage \section*{SEZNAM POUŽITÝCH ZKRATEK}
\addcontentsline{toc}{section}{SEZNAM POUŽITÝCH ZKRATEK}
\textit{abecedně seřadit\\pokud jsou anglicky, tak: český překlad (english meaning)}\\
\begin{tabular}{ll}
	CPL & kaprolaktam\\
	FID & plamenově-ionizační detektor pro GC\\
	GC  & plynová chromatografie\\
	LLE & rovnováha kapalina-kapalina (liquid-liquid equilibrium)\\
	TCE & trichlorethylen\\
	TOL & toluen\\
	TP  & teoretické patro\\
	SA  & síran amonný\\
\end{tabular}

\newpage \section*{SEZNAM SYMBOLŮ}
\addcontentsline{toc}{section}{SEZNAM SYMBOLŮ}
\textit{seřadit abecedně, a to: latinka (aAbBcC), řecká písmena, jiné znaky}

\singlespacing
\begin{tabular}{ll}
	$a$ & specifická mezifázová plocha [$\mathrm{m^{-1}}$]\\
	$A$ & plocha píku chromatogramu [$\mathrm{mV \cdot s}$]\\
	$A$ & průřez kolony [$\mathrm{m^2}$]\\
	$d_{32}$ & Sauterova střední velikost kapek [$\mathrm{m}$]\\
	$D$ & difuzivita kaprolaktamu v prostředí [$\mathrm{m^2 \cdot s^{-1}}$]\\
	$E$ & koeficient axiální disperze [$\mathrm{m^2 \cdot s^{-1}}$]\\
	$g$ & tíhové zrychlení [$\mathrm{m \cdot s^{-2}}$]\\
	$K$ & koeficient prostupu hmoty [$\mathrm{m \cdot s^{-1}}$]\\
	$m$ & hmotnost [$\mathrm{kg}$]\\
	$\dot{m}$ & hmotnostní tok [$\mathrm{kg \cdot s^{-1}}$]\\
	$Pe$ & Pécletovo číslo\\
	$Q$ & objemový tok [$\mathrm{m^3 \cdot s^{-1}}$]\\
	$r$ & efektivní poloměr kapiláry [$\mathrm{m}$]\\
	$S_{^i/_j}$ & kalibrační koeficient látky $i$ vůči $j$\\
	$T_x$ & počet přenosových jednotek vzhledem k fázi x\\
	$u$ & hmotnostní zlomek CPL v organické fázi\\
	$V$ & objem [$\mathrm{m^3}$]\\
	$w$ & hmotnostní zlomek CPL v organické fázi\\
	$\beta$ & koeficient přestupu hmoty [$\mathrm{m \cdot s^{-1}}$]\\
	$\gamma$ & mezifázové napětí [$\mathrm{N / m}$]\\
	$\rho$ & hustota [$\mathrm{kg \cdot m^{-3}}$]\\
	$\phi$ & zádrž dispergované fáze v kontinuální\\
	$\Phi$ & korekční faktor pro výpočet mezifázového napětí\\
\end{tabular}

\vspace{0.5cm} \textbf{Indexy}

\begin{tabular}{ll}
	* & hodnota v rovnováze\\
	C & kontinuální fáze\\
	D & dispergovaná fáze\\
	aq  & vodná fáze\\
	org & organická fáze\\
	x & fáze suroviny\\
	y & fáze rozpouštědla\\
	H & těžší fáze\\
	L & lehčí fáze\\
	F & surovina\\
	R & rafinát\\
	S & solvent\\
	E & extrakt\\
	in  & vstupní\\
	out & výstupní\\
\end{tabular}
\onehalfspacing

\newpage \section*{SEZNAM TABULEK}
\addcontentsline{toc}{section}{SEZNAM TABULEK}
\renewcommand{\listtablename}{\vskip -1cm}
\listoftables

\newpage \section*{SEZNAM OBRÁZKŮ}
\addcontentsline{toc}{section}{SEZNAM OBRÁZKŮ}
\renewcommand{\listfigurename}{\vskip -1cm}
\listoffigures

\newpage
\appendix
\section{ZDROJOVÝ KÓD} \label{sourceCode}
Nedílnou částí této práce je počítačový program, který je možné nalézt na přiloženém CD"-R nosiči.
Ten obsahuje zaprvé složku \textit{matlab}, v níž se nachází samotný program spustitelný v prostředí MATLAB (verze 2007 či novější), zadruhé soubor \textit{Dokumentace.pdf}, v němž je podrobně popsána struktura programu a použití jednotlivých funkcí včetně příkladů použití.

\end{document}
